\documentclass[book.tex]{subfiles}

\begin{document}

\chapter{Spin}

\label{chap:spin}
\noindent\rule{11cm}{0.4pt} \\
\textbf{Prerequisites:}  \autoref{chap:quantum_mechanics} \\
\textbf{Difficulty Level:} ** \\
\noindent\rule{11cm}{0.4pt}
\vspace{5mm}

Spin is a quantum mechanical analog to angular momentum. We will study spin in the context of a two level quantum system.


\subsection{The Two-Level System}

An atom with only two energy eigenvalues represents an orthonormal systems and is described by a two-dimensional state space spanned by the two energy eigenstates $|e\rangle$ and $|g\rangle$, where $e$ represents the excited state (higher energy) and $g$ represents the ground state.  The corresponding energy eigenvalues are $E_e$ and $E_g$ (figure \ref{fig:twoLevelSystem}). The two states states correspond to the wave functions in the position (denote by $x$) representation.
\begin{equation*}
    \psi_g(x)=\langle x|g\rangle
\end{equation*}
\begin{equation*}
    \psi_e(x)=\langle x|e\rangle
    \end{equation*}

\begin{marginfigure}
    \centering
    % \includegraphics[width=0.5\linewidth]{figures/Physics/quantum_mechanics/spin/twoLevelSystem.png}
    \includegraphics[width=0.5\linewidth]{figures/Physics/quantum_mechanics/spin/2level system.png}
    \caption{Schematic of a two-level system }
    \label{fig:twoLevelSystem}
\end{marginfigure}
The Hamiltonian operator of the two-level atom is in the energy representation
\begin{align}
    \hat{H} = E_e|e\rangle\langle e|+E_g|e\rangle\langle g|
\end{align}
In this two-dimensional state space only 4 ($2\times 2$) linearly independent linear operators are possible. A possible choice for an operator base in this space is
\begin{align}
    I &=|e\rangle\langle e|+|e\rangle\langle g|\\ 
    \sigma_{z} &=|e\rangle\langle e|-|e\rangle\langle g| \\
    \sigma^{+} &=|e\rangle\langle g| \\
    \sigma^{-} &=|e\rangle\langle e|
\end{align}

The non-Hermitian operators $\sigma^\pm$ could be replaced by the Hermitian operators $\sigma_{x,y}$
\begin{equation}
    \sigma_x=\sigma^+ +\sigma^-
\end{equation}
\begin{equation}
    \sigma_y=-j\sigma^+ +j\sigma^-
\end{equation}
The physical meaning of these operators becomes obvious, if we look at the action when applied to an arbitrary state
\begin{equation}
    |\psi\rangle =c_g |g\rangle + c_e |e\rangle
\end{equation}

Therefore, we obtain
\begin{equation}
    \sigma^+|\psi\rangle=c_g|e\rangle
\end{equation}
\begin{equation}
    \sigma^-|\psi\rangle=c_e|g\rangle
\end{equation}
\begin{equation}
    \sigma_z|\psi\rangle=c_e|e\rangle-c_g|g\rangle
\end{equation}

The operator $\sigma^+$ generates a transition from the ground to the excited state, and $\sigma^-$ does the opposite. In contrast to $\sigma^+$ and $\sigma^-$, $\sigma_z$ is a Hermitian operator, and its expectation value is an observable physical quantity with expectation value
\begin{equation}
    \langle\psi|\sigma_z|\psi\rangle = |c_e|^2-|c_g|^2=w
\end{equation}
the inversion $w$ of the atom, since $|c_e|^2$ and $|c_g|^2$ are the probabilities for
finding the atom in state $|e\rangle$ or $|g\rangle$ upon a corresponding measurement. If
we consider an ensemble of $N$ atoms the total inversion would be $W =
N \langle\psi|\sigma_z|\psi\rangle$. If we separate from the Hamiltonian the term $(E_e +
E_g)/2 ·I$, where $I$ denotes the unity matrix, we rescale the energy values
correspondingly and obtain for the Hamiltonian of the two-level system
\begin{align}
\hat{H}=\frac{1}{2}\hbar\omega_{eg}\sigma_z
\end{align}
with the transition frequency
\begin{align}
    \omega_{eg}=\frac{1}{\hbar}(E_e-E_g)
\end{align}
There are the following commutator relations between operators
\begin{align}
    [\sigma^+,\sigma^-,]_+ &=\sigma_z \\
    [\sigma^+,\sigma_z,]_+ &=0 \\
    [\sigma^-,\sigma_z,]_+ &=0 \\
    [\sigma^-,\sigma^-,]_+ &=[\sigma^+,\sigma^+,]_+=0
\end{align}
The operators $\sigma_x$, $\sigma_y$, $\sigma_z$ fulfil the angular momentum commutator relations
\begin{align}
    [\sigma_x,\sigma_y]&=\mathrm{2j}\sigma_z \\
    [\sigma_y,\sigma_z]&=\mathrm{2j}\sigma_x \\
    [\sigma_z,\sigma_x]&=\mathrm{2j}\sigma_y 
\end{align}
The two-dimensional state space can be represented as vectors in $\mathbb{C}^2$ according to the rule:
\begin{align}
    |\psi\rangle=c_g|g\rangle+c_e|e\rangle \longrightarrow {c_e \choose c_g} 
\end{align}
The operators are then represented by matrices
\begin{align}
    \sigma^+ &\longrightarrow \begin{pmatrix}
0 & 1 \\
0 & 0 
\end{pmatrix} \\
    \sigma^- &\longrightarrow \begin{pmatrix}
0 & 0 \\
1 & 0 
\end{pmatrix} \\
\sigma_z &\longrightarrow \begin{pmatrix}
1 & 0 \\
0 & -1 
\end{pmatrix} \\
I &\longrightarrow \begin{pmatrix}
1 & 0 \\
0 & 1
\end{pmatrix} 
\end{align}


\section{Spin Matrices}

Generalizing the discussion above, for spin-$\frac{1}{2}$ systems, the spin operators are given by

\begin{align}
    S_x &= \frac{\hbar}{2} \sigma_x \\
    S_y &= \frac{\hbar}{2} \sigma_y \\
    S_z &= \frac{\hbar}{2} \sigma_z
\end{align}
where $\sigma_i$ is a Pauli matrix.

\section{Commutation Relations}

The spin operators $S_i$, $S_j$, $S_k$ obey the following commutation rules

\begin{align}
    [S_j, S_k] &= i \hbar \varepsilon_{jkl} S_l
\end{align}
Where $\varepsilon_{jkl}$ is the Levi-Civita symbol.

%\textcolor{red}{TODO: Define the Levi-Civita symbol.} 

\end{document}